\section{Q\&A：Domain、Tokenizer、Fine-tuning 与 Dataset Release}

Q\&A 把 lecture recipe 放到边界场景：multimodal data 能否替代 text、code 是否需要特殊 architecture、一个 GPU 怎样做、scaling 是否跨 domain、tokenizer 有何差异、fine-tuning 是否用同样 filter，以及大数据集发布还需哪些治理。

\subsection{Code 与 text：相似 training，优先级不同}

讲者认为 architecture/training 总体相似，code model 仍可用 decoder Transformer；差别更多在 long repository context、FIM、fast IDE inference、MQA/GQA 和 data formatting。code tokenizer 仍用 BPE，但会保留 number splitting，并在 code mixture 上检查 under/overrepresented tokens。

\begin{knowledgebox}{Tokenizer 首次定义与 sanity checks}
BPE（Byte Pair Encoding）反复合并高频 token pairs。代码 tokenizer 应检查 whitespace/indentation、operators、numbers、Unicode、paths 和 language coverage；平均 token-per-character 过高会增加 context cost，过大的 whole-identifier tokens 又会损害组合泛化。
\end{knowledgebox}

\teachervoice{00:58:20--01:00:00，讲者说 code tokenizer 与 text tokenizer 很接近，重点是 number splitting、在真实 mixture 上训练并检查 outlier tokens；代码中也有大量 Markdown/English。}

\subsection{Pretraining 与 fine-tuning 的 filter 强度不同}

预训练需要 breadth，过强 star filter 会删掉太多数据；fine-tuning 面向特定 repository/language，可以更严格地筛高质量样本。LIMA 的 1,000 instructions 结果说明 alignment 阶段 quality 可能比 quantity 更重要，但不代表 1,000 样本适合所有任务。

\begin{warningbox}{不要把 LIMA 外推成“少数据总是够”}
LIMA 使用强 pretrained base 和特定 alignment setting。domain knowledge 若不在 base 中，fine-tuning 仍需要足够覆盖；小数据也更容易过拟合风格、泄漏评测或缺少 rare cases。
\end{warningbox}

\teachervoice{00:59:20--01:00:20，讲者对比：pretraining 不能轻易删到没有 breadth；fine-tuning 数据需求更小，可以投入更多精力做严格过滤。}

\subsection{Multimodal、Domain scaling 与不确定答案}

被问到 image/video 是否减少 text-only data 时，讲者明确说没有做过，不能给比例；她只观察到当时 VLM 仍有大量 text。被问到 medical 等 domain 时，她认为 scaling optimum 可能变化，现有 evidence 不足。

\begin{importantbox}{保留“不知道”也是 teacher voice}
高质量讲义不应替讲者补出确定答案。这里的可执行结论是：对新 modality/domain 重新做 scaling/data-mixture ablation，而不是引用 generic Chinchilla ratio。
\end{importantbox}

\teachervoice{00:52:40--00:55:00，讲者对 multimodal data substitution 直接回答“没有试过，不能真正回答”，只给出观察性判断。}

\subsection{发布大型数据集：技术、治理和访问控制一起设计}

发布前需要 license/copyright review、opt-out、PII/secrets detection、dataset card、processing code、statistics、gating sensitive subsets 和 versioned removal list。Hub/gated access 能改善可达性与风险控制，但不是法律豁免。

\begin{lstlisting}[caption={大型训练数据集发布前 checklist}]
record_source_license_timestamp_and_transform_lineage()
publish_filter_dedup_pii_and_decontamination_versions()
provide_inspection_opt_out_and_removal_workflows()
gate_sensitive_annotations_or_high_risk_subsets()
release_statistics_cards_tools_and_known_limitations()
version_updates_and_propagate_removals_to_future_models()
\end{lstlisting}

\teachervoice{01:00:00--01:01:30，讲者建议同时考虑 license、copyright、opt-out、工具和文档；对 PII annotation data 这类敏感 subset，可以使用 gated access。}

\subsection{本章小结}

Code LLM 并非完全不同的物种；它复用语言模型 recipe，却在 context、inference、format 和 evaluation 上有不同优先级。有限预算应做 curated fine-tuning；新 domain/modalities 则需要重新测 scaling；数据发布必须把治理工作流当成产品功能。

